\documentclass[11pt,a4paper]{article}

\usepackage[utf8]{inputenc}
\usepackage{amsmath,amssymb,amsthm}
\usepackage{geometry}
\usepackage{hyperref}
\usepackage{booktabs}
\usepackage{xcolor}

\geometry{margin=2.5cm}
\hypersetup{colorlinks=true,linkcolor=blue,citecolor=blue,urlcolor=blue}

\title{Markov Decision Process Specification\\for RL-Driven Physical Activity Interventions}
\author{William Dennis\\Bristol x NUS Research Project}
\date{June 2026}

\begin{document}

\maketitle

\section{High-Level Framework Design}

The \texttt{rl-health-interventions} framework simulates a wearable-device
physical activity intervention as a Markov Decision Process. Researchers
configure the data schema, environment dynamics, user behaviour, and agents
entirely through YAML config files --- no source code changes required.

\subsection{Architecture}

The framework is organised into five subphases, each delivering a pluggable
component:

\begin{enumerate}
    \item \textbf{Data Layer (1A):} Ingests CSV/Parquet files via Polars lazy
    scans, applies config-driven feature engineering, and materialises a
    numpy-backed \texttt{Dataset}. Provides a \texttt{StateView} bridge for
    the simulation loop.
    \item \textbf{MDP Environment (1B):} Exposes a clean
    \texttt{step(state, action) $\rightarrow$ (next\_state, reward, done)} API.
    Pluggable \texttt{TransitionModel} and \texttt{RewardHandler} components
    are selected from a registry by config string.
    \item \textbf{User Simulation (1C):} Rule-based behavioural archetypes
    (goal-driven, social responder, resistant, stable maintainer) with
    backlash and fatigue mechanisms. A \texttt{ResponseModel} maps
    $(\text{state}, \text{action}, \text{profile})$ to a next state.
    \item \textbf{Agent Library (1D):} Common \texttt{Agent} ABC with
    \texttt{select\_action(state) $\rightarrow$ int} and
    \texttt{update(...)}. Thompson Sampling is the MVP gate; deep RL
    agents (DQN, PPO, CQL, IQL) are stretch.
    \item \textbf{Experiment Runner (1E):} Wires all components via a single
    \texttt{ExperimentFactory.build(config)}. Runs multi-trial, multi-agent
    experiments with config snapshot reproducibility.
\end{enumerate}

\subsection{Component Registration}

Each pluggable family (transitions, rewards, agents, response models) is a flat
package with a \texttt{\_base.py} ABC and a \texttt{REGISTRY} dict in its
\texttt{\_\_init\_\_.py}. Adding a new component requires one file, one registry
entry, and a config string reference --- no framework code changes.

\subsection{Dependency Graph}

Tracks A (1A) and B (1B) run in parallel. Track C (Dataset Exploration)
runs alongside both and is now complete (\texttt{docs/03 Data Sources.md}).
After 1A, 1B, and exploration complete, tracks merge at
1C (user simulation) and 1D (agent library, which only needs the 1B interface).
The final track is 1E (experiment runner), which depends on 1C and 1D.

\subsection{Validation}

The factory applies three layers of validation before the main experiment loop:
(1) Pydantic schema parsing rejects malformed configs at load time; (2) component
compatibility checks verify registry keys and cross-references; (3) a single
dummy \texttt{step()} catches wiring errors and shape mismatches.

\subsection{Test Organisation}

Tests mirror the component module structure (\texttt{tests/unit/transitions/},
\texttt{tests/unit/rewards/}, etc.), not subphase numbers. Every ABC method has
at least one test. Integration and end-to-end tests cover factory wiring and the
full experiment pipeline.

\section{MDP Specification}

This document specifies the Markov Decision Process (MDP) for the first iteration
of the \texttt{rl-health-interventions} simulation framework. It formalises the
design from the project google doc and incorporates decisions from the
2026-06-08 planning session.

The MDP is defined as the tuple $(\mathcal{S}, \mathcal{A}, P, R, \gamma)$.

\section{State Space $\mathcal{S}$}

The state represents the current context of a user at decision epoch $t$.
States are constructed from wearable-device signals, user profile information,
and intervention history:

\[
s_t = \{ \mathbf{x}_t^{\text{wearable}}, \mathbf{x}_t^{\text{context}},
\mathbf{x}_t^{\text{history}}, \mathbf{x}^{\text{profile}} \}
\]

\begin{table}[h]
\centering
\caption{State variables}
\begin{tabular}{lll}
\toprule
Variable & Type & Description \\
\midrule
$\text{steps}_t$ & $\mathbb{R}_{\geq 0}$ & Step count, past 24h \\
$\text{hr}_t$ & $\mathbb{R}_{\geq 0}$ & Heart rate (or resting) \\
$\text{sleep\_hours}_t$ & $\mathbb{R}_{\geq 0}$ & Previous night sleep \\
$\text{sedentary\_min}_t$ & $\mathbb{R}_{\geq 0}$ & Sedentary minutes today \\
$\text{time\_of\_day}_t$ & $\{0,\dots,23\}$ & Current hour \\
$\text{day\_of\_week}_t$ & $\{0,\dots,6\}$ & Day of week \\
$\text{goal\_progress}_t$ & $[0,1]$ & Fraction of daily goal met \\
$\text{burden}_t$ & $\mathbb{Z}_{\geq 0}$ & Interventions sent today \\
$a_{t-1}$ & $\mathcal{A}$ & Previous action \\
$\text{response}_{t-1}$ & $\{0,1\}$ & Previous response \\
$\text{body\_measure}_k$ & $\mathbb{R}$ & Clinical measure (q3wks) \\
$\text{age}$ & $\mathbb{Z}_{>0}$ & Static profile \\
$\text{gender}$ & $\{0,1\}$ & Static profile \\
$\text{baseline\_activity}$ & $\{0,1,2\}$ & Low / medium / high \\
\bottomrule
\end{tabular}
\end{table}

Variables are normalised to $[0,1]$ before being passed to the agent. The body
measure is only observed every 3 weeks (sparse reward signal).

\section{Action Space $\mathcal{A}$}

\[
\mathcal{A} = \{ a_0, a_1, a_2, a_3, a_4, a_5 \}
\]

\begin{table}[h]
\centering
\caption{Action space}
\begin{tabular}{cll}
\toprule
Action & Label & Description \\
\midrule
$a_0$ & No message & Do not intervene \\
$a_1$ & Motivational prompt & Generic encouragement \\
$a_2$ & Walking suggestion & Recommend a short walk \\
$a_3$ & Goal reminder & Remind of step goal \\
$a_4$ & Recovery suggestion & Stretch or rest \\
$a_5$ & Progress feedback & Report goal progress \\
\bottomrule
\end{tabular}
\end{table}

Exactly one action per decision epoch. The ``no message'' action is critical
for managing notification burden.

\section{Transition Function $P$}

The transition function models the probability over next states:

\[
P(s_{t+1} \mid s_t, a_t): \mathcal{S} \times \mathcal{A} \times \mathcal{S}
\rightarrow [0,1]
\]

In Phase 1 (rule-based), transitions follow a behavioural response model:

\[
\Delta \text{steps}_{t+1} = f(s_t, a_t, \theta_{\text{user}})
\]

where $\theta_{\text{user}}$ parameterises one of four behavioural archetypes:

\begin{itemize}
    \item \textbf{Goal-driven:} Responds to reminders and feedback.
    Proportional to goal progress.
    \item \textbf{Social responder:} Responds to motivational prompts.
    Models social desirability bias.
    \item \textbf{Resistant:} Low response. Fatigue accumulates quickly.
    \item \textbf{Stable maintainer:} Already active. Low marginal gain.
\end{itemize}

Burden evolves as:

\[
\text{burden}_{t+1} =
\begin{cases}
\text{burden}_t + 1 & \text{if } a_t \neq a_0 \\
0 & \text{if } a_t = a_0
\end{cases}
\]

When $\text{burden} > B_{\max}$, response probability decays linearly.

\section{Reward Function $R$}
\label{sec:reward}

\subsection{Immediate Reward}

\[
R_t^{\text{immediate}} = \alpha \cdot \Delta\text{steps}_t
- \beta \cdot \mathbf{1}_{\{a_t \neq a_0\}}
+ \lambda \cdot \text{goal\_progress}_t
\]

\subsection{Delayed Reward (Body Measures)}

Clinical body measures arrive every 3 weeks ($t = 3k$):

\[
R_{3k}^{\text{delayed}} = \eta \cdot \text{BM\_improvement}_{3k}
\]

Total reward:

\[
R_t =
\begin{cases}
R_t^{\text{immediate}} + R_t^{\text{delayed}} & t \equiv 0 \pmod{3} \\
R_t^{\text{immediate}} & \text{otherwise}
\end{cases}
\]

The compound reward (immediate + delayed) is a central research question.
Weights $(\alpha, \beta, \lambda, \eta)$ are configurable experiment parameters.

\section{Discount Factor $\gamma$}

\[
\gamma \in [0.9, 0.95]
\]

A high discount factor encourages long-term behaviour change over short-term
step increases. Appropriate for health habit formation.

\section{Decision Epoch}

\textbf{Daily} decisions. This aligns with the weekly intervention cadence,
the 3-week body measure delay (3 epochs), and practical deployment constraints.

\section{Configurability}

All MDP components are exposed in the config schema: state variable selection,
action subsets, reward weights, transition archetypes, fatigue threshold,
discount factor, and epoch frequency. New datasets or intervention designs
require a new config file, not new code.

\end{document}
